% ============================================================================
%  ECR methodology - two-page mathematical breakdown.
%  Follows on from \section{ECR Data-Driven Targeting}\label{sec:ecr-method}.
%  Requires amsmath, amssymb, bm (and amsthm if the definitions are kept as
%  theorem environments; they are plain text here to save space).
%  The longer treatment, with the local geometry of the target, the training
%  algorithm and the remarks on delay coordinates, is in
%  ecr_method_math_full.tex.
%
%  NOTATION: the observed (controller) state is bold, \mathbf{z}_n; plain
%  x, y, z are the Lorenz coordinates.  Rename to \bm{\zeta}_n if the clash
%  with the Lorenz z is unwelcome.
% ============================================================================

\subsection{Plant, surface of section and admissible perturbations}
\label{subsec:ecr-plant}

The Lorenz system is written in full, since the second equation defines the
surface of section used throughout,
%
\begin{equation}
\dot{x}=\sigma\,(y-x),
\qquad
\dot{y}=\rho\,x-y-xz,
\qquad
\dot{z}=xy-\beta z,
\label{eq:lorenz}
\end{equation}
%
with state $\mathbf{u}=[x\;y\;z]^{\mathsf T}$, flow
$\varphi_t(\mathbf{u};\mathbf{p})$ and control parameter vector
$\mathbf{p}=[\sigma\;\rho\;\beta]^{\mathsf T}\in\mathbb{R}^{r}$, $r=3$. The
section is the half-plane
%
\begin{equation}
\Sigma=\bigl\{\mathbf{u}:y=y_{\mathrm{PSS}},\;\dot y<0\bigr\},
\qquad
y_{\mathrm{PSS}}=\sqrt{\beta_{\mathrm{nom}}(\rho_{\mathrm{nom}}-1)}=8.4853,
\label{eq:section}
\end{equation}
%
transversality requiring $\dot y|_{\Sigma}=\rho x-y_{\mathrm{PSS}}-xz\neq0$ by
\eqref{eq:lorenz}; the sign condition picks the branch that reproduces the
target of the source paper, the opposite one containing the equilibrium
$C^{+}$ itself. With the first return time
$\tau(\mathbf{u};\mathbf{p})=\min\{t>0:\varphi_t(\mathbf{u};\mathbf{p})\in\Sigma\}$
and $y\equiv y_{\mathrm{PSS}}$ on $\Sigma$, the observed state at the $n$-th
crossing is $\mathbf{z}_n=[x(t_n)\;z(t_n)]^{\mathsf T}\in\mathbb{R}^{N}$,
$N=2$, and the dynamics reduce to the discrete map
%
\begin{equation}
\mathbf{z}_{n+1}=\mathbf{G}(\mathbf{z}_n,\mathbf{p}_n),
\label{eq:map}
\end{equation}
%
the parameter being held constant over the flight $t\in[t_n,t_{n+1})$ and
changed only at crossings. If only a scalar measurement $s(t)$ is available, the
delay vector $[s(t_n)\;s(t_n-T)\;\cdots]^{\mathsf T}$ replaces the section state
and \eqref{eq:map} still holds generically. No closed form of $\mathbf{G}$ is
ever constructed.

Each control parameter is confined to
%
\begin{equation}
\Pi=\Bigl\{\mathbf{p}:\bigl|p^{i}-p^{i}_{\mathrm{nom}}\bigr|<\delta p^{i}_{\max},\;i=1,\dots,r\Bigr\},
\qquad
\Delta\mathbf{p}_n=\bigl(\mathbf{p}_n-\mathbf{p}_{\mathrm{nom}}\bigr)\oslash\delta\mathbf{p}_{\max},
\label{eq:admissible-set}
\end{equation}
%
with $\mathbf{p}_{\mathrm{nom}}=[10\;28\;8/3]^{\mathsf T}$ and
$\delta\mathbf{p}_{\max}=[0.30\;0.84\;0.08]^{\mathsf T}$, i.e.\ $3\%$ of each
nominal value --- a design choice carried over from the source paper rather
than a property of the method, $\Pi$ needing only to be small enough to leave
the attractor intact --- so that admissibility is the box test
$\mathbf{p}_n\in\Pi\iff\|\Delta\mathbf{p}_n\|_{\infty}\le1$ on the normalised
perturbation $\Delta\mathbf{p}_n$. The target is a fixed point of the controlled
map at nominal parameters,
$\mathbf{z}^{*}=\mathbf{G}(\mathbf{z}^{*},\mathbf{p}_{\mathrm{nom}})$, i.e.\ a
period-1 unstable periodic orbit of the flow, located by Newton
iteration on $\mathbf{G}(\mathbf{z},\mathbf{p}_{\mathrm{nom}})-\mathbf{z}$ with a
numerically differentiated Jacobian: here
$\mathbf{z}^{*}=[14.2522\;\,39.7867]^{\mathsf T}$, with map eigenvalues
$\lambda_1=4.71$ and $\lambda_2\approx10^{-9}$, a saddle so strongly contracting
that the return map is effectively one-dimensional. The sensitivity
$\partial\mathbf{G}/\partial\mathbf{p}$ there has full rank $N$, and the
displacement reachable within $\Pi$ exceeds the target radius $\delta=0.30$ by
an order of magnitude, which is what makes single-step targeting feasible.

\subsection{Control regions and their empirical construction}
\label{subsec:ecr-regions}

With $\mathcal{B}_\delta=\{\mathbf{z}:\|\mathbf{z}-\mathbf{z}^{*}\|<\delta\}$,
the target region and the extended control regions are
%
\begin{align}
T_0&=\Bigl\{\mathbf{z}\in\mathcal{B}_\delta:\exists\,\mathbf{p}\in\Pi,\;
\mathbf{G}(\mathbf{z},\mathbf{p})\in\mathcal{B}_\delta\Bigr\},
\label{eq:def-T0}\\[2pt]
T_i&=\Bigl\{\mathbf{z}\notin\textstyle\bigcup_{j<i}T_j:\exists\,\mathbf{p}\in\Pi,\;
\mathbf{G}(\mathbf{z},\mathbf{p})\in T_{i-1}\Bigr\},
\qquad i=1,\dots,K .
\label{eq:def-Ti}
\end{align}
%
The regions are disjoint by construction and membership of $T_i$ certifies
capture of $T_0$ within $i$ admissible steps; their union is the activation
region, outside which $\mathbf{p}_{\mathrm{nom}}$ is applied. Both are existence
statements about admissible perturbations, not formulae in $\mathbf{G}$.

The plant is run from many initial conditions while an independent
$\mathbf{p}_n\sim\mathcal{U}(\Pi)$ is applied at each crossing, giving the
transition set
$\mathcal{D}=\{(\mathbf{z}_n,\mathbf{p}_n,\mathbf{z}_{n+1})\}_{n=1}^{M}$, in
which each existential quantifier above is settled by observed evidence, the
applied $\mathbf{p}_n$ being the witness:
%
\begin{align}
\mathcal{I}_0&=\bigl\{n:\|\mathbf{z}_n-\mathbf{z}^{*}\|<\delta\;\wedge\;
\|\mathbf{z}_{n+1}-\mathbf{z}^{*}\|<\delta\bigr\},
\label{eq:index-T0}\\[2pt]
\mathcal{I}_i&=\bigl\{n:\mathcal{M}_{i-1}(\mathbf{z}_{n+1})=1\;\wedge\;
\mathcal{M}_{j}(\mathbf{z}_{n})=0\;\;\forall j<i\bigr\},
\qquad i\ge1,
\label{eq:index-Ti}
\end{align}
%
membership of the successor in $T_{i-1}$ being decided by the already-trained
level-$(i-1)$ model acting as an oracle
$\mathcal{M}_{i-1}:\mathbb{R}^{N}\to\{0,1\}$, defined in
\eqref{eq:membership-oracle}. The training pairs of region $i$ are
$\{(\mathbf{z}_n,\Delta\mathbf{p}_n)\}_{n\in\mathcal{I}_i}$ --- the state, and the
perturbation observed to move it one region inward --- and construction
proceeds bottom-up, each level supplying the oracle for the next.

\subsection{Clustering, region assignment and local models}
\label{subsec:ecr-clustering}

ECR-I fits one network over the whole activation region, ECR-II one per region,
and ECR-III partitions each region first, since data in the outer regions is
scattered. On $\mathcal{Z}_i=\{\mathbf{z}_n\}_{n\in\mathcal{I}_i}$ the clusters
$C_{i1},\dots,C_{in_i}$ are the equivalence classes of the transitive closure
of $\mathbf{z}\sim\mathbf{z}'\iff\|\mathbf{z}-\mathbf{z}'\|\le r_i$, with $r_i$
taken at the first minimum of the histogram of pairwise distances. Each cluster
carries its mean $\bm{\mu}_{ij}$ and a ridge-regularised covariance
$\bm{\Sigma}_{ij}$, and a state is assigned to the cluster minimising the
normalised Mahalanobis distance
%
\begin{equation}
\mathrm{NMD}^{2}_{ij}(\mathbf{z}_n)=(\mathbf{z}_n-\bm{\mu}_{ij})^{\mathsf T}\,
\underline{\mathbf{N}}_{ij}\,(\mathbf{z}_n-\bm{\mu}_{ij}),
\qquad
\underline{\mathbf{N}}_{ij}=\Bigl(\bm{\Sigma}_{ij}\big/\det(\bm{\Sigma}_{ij})^{1/N}\Bigr)^{-1},
\label{eq:nmd}
\end{equation}
%
the unit-determinant normalisation removing the size of a cluster but not its
shape, so clusters of different populations compete on equal terms (plain,
variance- and trace-normalised readings are the alternatives). Each cluster
carries a Gaussian RBF network mapping the standardised state
$\tilde{\mathbf{z}}$ to a normalised perturbation,
%
\begin{equation}
\widehat{\Delta\mathbf{p}}(\mathbf{z})=\mathbf{W}^{\mathsf T}
\bigl[\phi_1\;\cdots\;\phi_{n_c}\;1\bigr]^{\mathsf T},
\quad
\phi_k=\exp\!\Bigl(-\tfrac{\|\tilde{\mathbf{z}}-\mathbf{c}_k\|^{2}}{2s_k^{2}}\Bigr),
\quad
\mathbf{W}=\bigl(\bm{\Phi}^{\mathsf T}\bm{\Phi}+\lambda\mathbf{I}\bigr)^{-1}\bm{\Phi}^{\mathsf T}\mathbf{Y},
\label{eq:rbf}
\end{equation}
%
with centres from $k$-means, widths from the spacing of neighbouring centres,
and a closed-form solve in place of iterative training. The oracle of
\eqref{eq:index-Ti} is then an admissibility test on the nearest cluster's
network,
%
\begin{equation}
\mathcal{M}_i(\mathbf{z})=\mathbb{1}\Bigl[\bigl\|\widehat{\Delta\mathbf{p}}^{\,(ij^{*})}(\mathbf{z})\bigr\|_{\infty}\le1\Bigr],
\qquad
j^{*}=\arg\min_j\mathrm{NMD}_{ij}(\mathbf{z}),
\label{eq:membership-oracle}
\end{equation}
%
optionally gated by a multiple of the $95$th percentile of the cluster's own
training NMD, since \eqref{eq:nmd} returns a nearest cluster even for states
far outside every recorded patch.

\subsection{Control law and performance measures}
\label{subsec:ecr-online}

At each crossing the state is measured, a region is identified and the
corresponding perturbation applied:
%
\begin{equation}
\mathbf{p}_n=
\begin{cases}
\mathbf{p}_{\mathrm{nom}}+\widehat{\Delta\mathbf{p}}^{\,(i^{*}j^{*})}(\mathbf{z}_n)\odot\delta\mathbf{p}_{\max},
& \text{if an admissible } (i^{*},j^{*}) \text{ exists},\\
\mathbf{p}_{\mathrm{nom}}, & \text{otherwise},
\end{cases}
\label{eq:control-law}
\end{equation}
%
so the applied perturbation is admissible by construction. The variants differ
only in how $(i^{*},j^{*})$ is found: ECR-II searches the regions in ascending
order and accepts the first admissible answer,
$i^{*}=\min\{i:\|\widehat{\Delta\mathbf{p}}^{\,(i)}(\mathbf{z}_n)\|_{\infty}\le1\}$,
whereas ECR-III obtains it analytically as
$(i^{*},j^{*})=\arg\min_{i,j}\mathrm{NMD}_{ij}(\mathbf{z}_n)$ over the clusters of
all regions at once, trading a sequential search over networks for a distance
computation. Performance is measured by the reaching time and its ensemble mean
over initial conditions on the attractor,
%
\begin{equation}
n_r(\mathbf{z}_0)=\min\bigl\{n\ge0:\|\mathbf{z}_n-\mathbf{z}^{*}\|<\delta\bigr\},
\qquad
\bar{n}_r=\frac{1}{M_{\mathrm{ic}}}\sum_{m=1}^{M_{\mathrm{ic}}}
\min\bigl\{n_r(\mathbf{z}_0^{(m)}),n_{\max}\bigr\},
\label{eq:reaching-time}
\end{equation}
%
censored at a horizon $n_{\max}$, together with the capture rate
$\Pr\{n_r<n_{\max}\}$ and the retention, the fraction of steps after capture
spent inside $\mathcal{B}_\delta$, which separates successful targeting from
successful stabilisation. Without targeting $\bar{n}_r$ is the natural
recurrence time to $\mathcal{B}_\delta$, some $15$ steps here.
